% =====================================================================
% ANTECEDENTES -- Responsable: Fabian Moreno Ugarte
%
% Fuente: los tres PDF de ./papers/.
%
% CRITERIO APLICADO A LAS CIFRAS (no modificar sin leer NOTAS_FUENTES.md):
% solo se reproducen las cifras que constan en el TEXTO EN PROSA de cada
% articulo (resumen, introduccion o discusion de resultados). Las que
% figuraban unicamente en las tablas de resultados fueron retiradas,
% porque la conversion de esos cuadros a texto plano mezcla filas y
% columnas y una cifra mal atribuida a un conjunto de datos seria
% exactamente el tipo de afirmacion no verificable que este informe debe
% evitar. La trazabilidad esta en NOTAS_FUENTES.md, seccion "Papers".
% =====================================================================
% 14/09/2026: pasa a ser la Seccion 4 (Marco Teorico) del indice del
% modelo del curso. Contenido sin cambios; se conserva sec:antecedentes.
\section{Marco Teórico}
\label{sec:antecedentes}

El conteo de personas en escenas congestionadas es un problema con una
línea de investigación propia dentro de la visión computacional. Esta
sección revisa tres trabajos que marcan tres etapas sucesivas de esa
línea (la regresión de mapas de densidad, la detección puramente
basada en puntos y el aprendizaje semi-supervisado sobre puntos) y
discute, para cada uno, qué propone, por qué resulta pertinente para el
caso del Aeropuerto Internacional Jorge Chávez y qué limitación
reconocen sus propios autores.

Se privilegia deliberadamente la limitación \emph{declarada por los
autores} sobre la crítica externa: en un informe de cumplimiento, lo que
un proveedor de tecnología reconoce por escrito sobre su propio método
es el insumo más defendible para dimensionar el riesgo residual del
sistema.

Solo se reproducen aquí las cifras que constan en el texto en prosa de
cada artículo, es decir, en su resumen, introducción o discusión de
resultados. Las que
figuraban únicamente en las tablas de resultados se omiten: la conversión
de esos cuadros a texto plano mezcla filas y columnas, y una cifra mal
atribuida a un conjunto de datos sería precisamente el tipo de afirmación
no verificable que este informe debe evitar. Donde la comparación importa,
se citan las magnitudes relativas que los propios autores enuncian en
prosa, que es la forma en que además presentan sus contribuciones.

% ---------------------------------------------------------------------
\subsection{CSRNet: redes convolucionales dilatadas para escenas altamente congestionadas}
\label{sec:ant-csrnet}

\subsubsection*{Qué propone}

Li, Zhang y Chen~\cite{Li_2018_CVPR} proponen CSRNet, una red
completamente convolucional para reconocimiento de escenas congestionadas
compuesta por dos bloques: un \emph{front-end} que reutiliza las primeras
diez capas convolucionales de VGG-16 para extracción de características
2D, y un \emph{back-end} de convoluciones dilatadas que amplía el campo
receptivo sin recurrir a capas de \emph{pooling}. La motivación explícita
es prescindir de las arquitecturas multi-columna previas: los autores
muestran experimentalmente que las tres columnas de MCNN aprenden
características casi idénticas, de modo que la ramificación introduce
redundancia y tiempo de entrenamiento sin aportar diversidad de escala
real. La salida del modelo no es un número sino un mapa de densidad,
entrenado con pérdida euclídea contra un mapa de referencia generado
difuminando cada anotación de cabeza con un núcleo gaussiano
geométricamente adaptativo.

Los autores enuncian sus resultados en términos comparativos frente al
estado del arte del momento: sobre ShanghaiTech obtienen un 7\,\% menos de
MAE que CP-CNN en la Parte A y un 47,3\,\% menos en la Parte B. Extienden
además el método al conteo de vehículos sobre TRANCOS, donde reportan un
15,4\,\% menos de MAE que el mejor método previo.

\subsubsection*{Por qué aplica al caso LAP}

Tres propiedades del diseño son directamente relevantes para el terminal.

\begin{enumerate}[leftmargin=1.6em,itemsep=0.3em]
  \item \textbf{La salida es un mapa de densidad, no una identidad.} El
        modelo produce una superficie continua de concentración espacial
        y el conteo se obtiene integrándola. Ninguna etapa del
        \emph{pipeline} requiere extraer rasgos faciales ni construir una
        plantilla biométrica. Como se argumenta en la
        Sección~\ref{sec:auditoria}, esta característica es el principal
        insumo técnico para sostener que el tratamiento no recae sobre
        datos biométricos, aunque la conclusión depende también de qué se
        almacene aguas arriba del modelo.
  \item \textbf{La arquitectura es puramente convolucional y admite
        resoluciones variables.} Esto importa operativamente: las cámaras
        de un terminal no comparten resolución, altura ni ángulo, y el
        modelo no obliga a reescalar todo a un formato único.
  \item \textbf{La distribución espacial es informativa por sí misma.}
        Los autores ilustran el punto con tres imágenes que contienen el
        mismo número de personas pero distribuciones radicalmente
        distintas. Para LAP, saber que hay 300 personas en el hall de
        salidas es mucho menos útil que saber que 200 de ellas están
        concentradas frente a tres mostradores.
\end{enumerate}

\subsubsection*{Limitación reconocida por los autores}

La limitación más explícita aparece en la evaluación sobre UCSD, un
conjunto de escenas relativamente dispersas (entre 11 y 46 personas por
imagen): los propios autores indican que superan a la mayoría de métodos
previos \emph{excepto a MCNN en la categoría MAE}. El artículo consigna
en prosa que CSRNet alcanza 1,16 de MAE sobre ese conjunto, pero no es
esa cifra sino la excepción reconocida la que interesa aquí. Reconocen
además que la baja resolución de
los cuadros ($238\times158$ píxeles) dificulta generar un mapa de
densidad de calidad tras operaciones de \emph{pooling} sucesivas, y que
fue necesario reescalar por interpolación bilineal a $952\times632$.

Esto tiene una consecuencia directa para el proyecto: el régimen de baja
densidad no es el terreno donde el método brilla, y buena parte del día
un terminal aéreo opera precisamente en baja densidad. El método está
optimizado para el escenario que menos ocurre.

Una segunda limitación, más estructural, es que el mapa de densidad no
entrega posiciones individuales. Los autores lo asumen como parte del
paradigma; el trabajo siguiente lo convierte en su crítica central.

% ---------------------------------------------------------------------
\subsection{P2PNet: un marco puramente basado en puntos}
\label{sec:ant-p2pnet}

\subsubsection*{Qué propone}

Song et al.~\cite{Song_2021_ICCV} sostienen que localizar individuos
responde mejor a las demandas prácticas del análisis de multitudes que
simplemente contarlos, y que los métodos existentes basados en
representaciones intermedias (mapas de densidad o cajas delimitadoras
pseudo-generadas) son contraintuitivos y propensos a error. Su
propuesta, P2PNet, predice directamente un conjunto de puntos con sus
coordenadas y confianzas, sin pasar por ninguna representación
intermedia. La pieza crítica del método es la asignación de objetivos de
aprendizaje: los autores establecen que tanto el caso en que varias
propuestas se emparejan con una misma anotación como el caso inverso
confunden al modelo durante el entrenamiento y producen conteos sobre- o
sub-estimados, y resuelven el problema con un emparejamiento uno a uno
mediante el algoritmo húngaro.

El trabajo aporta además una métrica, la \emph{density Normalized Average
Precision} (nAP), que evalúa conjuntamente localización y conteo, y que
fue diseñada para no ignorar la variación de densidad entre escenas ni
dejar sin penalización las predicciones duplicadas.

En cuanto a resultados de conteo, los autores los describen en prosa de
forma comparativa: sobre ShanghaiTech Parte A reducen el MAE un 4,8\,\% y
el MSE un 12,9\,\% respecto del segundo mejor método, ADSCNet, y sobre la
Parte B logran una reducción del 2,3\,\% en MAE; sobre UCF\_CC\_50 rebajan
el MAE en 2,1 puntos frente al segundo mejor resultado, y sobre
NWPU-Crowd un 12,4\,\% frente a DM-Count. La única cifra absoluta que el
artículo enuncia fuera de sus tablas es un MAE de 85,32 sobre UCF-QNRF.
En localización, los autores señalan que el nAP$_{0.5}$ se sitúa cerca del
90\,\% en la mayoría de los conjuntos y reportan un F1/Precisión/Recall de
71,2\,\%/72,9\,\%/69,5\,\% bajo la métrica de NWPU-Crowd.

\subsubsection*{Por qué aplica al caso LAP}

La localización individual es lo que convierte una cifra en una decisión
operativa. Un conteo agregado por cámara no permite decidir dónde abrir
un mostrador adicional o hacia dónde redirigir el flujo; un conjunto de
puntos sí, porque admite agregación por zonas definidas por LAP (colas de
\emph{check-in}, filtro de seguridad, migraciones) sin reentrenar el
modelo.

Ahora bien, esta misma capacidad es la que eleva el perfil de riesgo del
sistema y obliga a un análisis más fino en la Sección~\ref{sec:marco}. Un
punto con coordenadas no es un dato biométrico, pero sí es un dato
espacio-temporal referido a una persona física presente en un lugar y
momento determinados. Si esos puntos se persisten y se encadenan en el
tiempo, el sistema se acerca a un seguimiento de trayectorias, que es una
finalidad distinta de la estimación de aglomeraciones y que requeriría su
propia base de legitimación.

Queda así un riesgo a validar con el cliente. La diferencia entre
\emph{estimar una aglomeración} y \emph{seguir a una persona} no es una
diferencia de algoritmo sino de política de persistencia. Con la misma
salida de P2PNet, retener los puntos y asociarlos entre cuadros produce
trayectorias individuales. Este límite debe fijarse por diseño y por
contrato, no dejarse a la configuración: una vez que el sistema está en
producción, permitir la persistencia es un cambio de parámetro, mientras
que prohibirla requiere volver sobre el diseño.

\subsubsection*{Limitación reconocida por los autores}

Los autores reconocen tres limitaciones de manera explícita.

Primero, en UCF-QNRF admiten que su precisión \enquote{no es tan
competitiva} frente a ADSCNet, aunque sea superior en todos los demás
conjuntos. Segundo, en NWPU-Crowd atribuyen su desempeño inferior en la
métrica MAE[S] a que sus predicciones se basan en un único mapa de
características de una sola escala, elegido por simplicidad; señalan que
incorporar fusión multiescala sería una mejora natural. Tercero, y más
relevante para el caso LAP, observan que la precisión es notablemente
menor bajo el umbral más estricto nAP$_{0.05}$, y explican que a esa
exigencia de localización los efectos de las desviaciones de etiquetado
empiezan a hacerse aparentes. El artículo no cuantifica esa caída fuera
de sus tablas; el contraste que sí enuncia en prosa es que el
nAP$_{0.5}$ ronda el 90\,\% mientras que ya bajo el umbral intermedio
nAP$_{0.25}$ la precisión se mantiene por encima del 55\,\%.

Esta última observación es un límite \emph{del suelo de verdad}, no del
modelo: por debajo de cierta escala, la propia anotación humana es
ruidosa. Cualquier compromiso de exactitud que el equipo ofrezca a LAP
debe formularse a nivel de zona agregada y no a nivel de posición
individual, porque a nivel individual la referencia contra la cual se
mediría no es confiable.

% ---------------------------------------------------------------------
\subsection{P2R: pérdida punto-a-región para conteo semi-supervisado}
\label{sec:ant-p2r}

\subsubsection*{Qué propone}

Lin, Zhao y Chan~\cite{Lin_2025_CVPR} parten de una constatación
económica: entrenar un contador basado en puntos exige anotar cientos o
miles de puntos por imagen, y ese costo ha frenado la adopción del
enfoque. Integran entonces los métodos basados en puntos en un marco
semi-supervisado de pseudo-etiquetado, en el que un modelo maestro genera
etiquetas para imágenes no anotadas y un modelo estudiante aprende de
ellas.

Al implementarlo detectan un problema específico: la confianza asociada a
una pseudo-etiqueta no logra propagarse a los píxeles de fondo bajo el
esquema punto-a-punto. Para diagnosticarlo diseñan un mapa de activación
específico por punto (PSAM), y observan que los píxeles vecinos de cada
píxel de primer plano quedan sobre-activados, de modo que el
decodificador los interpreta como instancias distintas. La solución
propuesta es sustituir el emparejamiento punto-a-punto por uno
punto-a-región (P2R): en lugar de detectar un punto por peatón, se
segmenta una región local, de forma que todos los píxeles de esa región
comparten la confianza del pseudo-punto correspondiente. Como efecto
colateral, el algoritmo húngaro deja de ser necesario, lo que reduce el
costo de cómputo de la pérdida.

Los experimentos se realizan sobre cuatro conjuntos (ShanghaiTech A y
B, UCF-QNRF y JHU-Crowd++) bajo protocolos de 5\,\%, 10\,\% y 40\,\%
de datos etiquetados. La afirmación central que los autores enuncian en
prosa es que entrenar con el 5\,\% de los datos usando P2R supera al
aprendizaje totalmente supervisado con el 10\,\%, y resulta equivalente a
otros métodos semi-supervisados que emplean ese mismo 10\,\%. Como
referencia del punto de partida, el artículo consigna que el modelo
entrenado únicamente con ese 5\,\% de datos etiquetados alcanza un MAE de
93,7 y un MSE de 155,2 sobre ShanghaiTech Parte A. Bajo supervisión
completa (100\,\% de etiquetas) P2R supera a P2PNet en los cuatro
conjuntos.

La supresión del algoritmo húngaro tiene además un efecto que los autores
miden y describen en prosa: sobre una imagen de $576\times960$ con 775
puntos anotados, el cálculo de la pérdida punto-a-punto requiere 0,4307
segundos frente a 0,0064 segundos de P2R, cerca de 68 veces más rápido.
No es un detalle menor para el proyecto, porque el costo de cómputo del
entrenamiento tiene una lectura ambiental que se retoma en la
Sección~\ref{sec:impacto-ambiental}.

\subsubsection*{Por qué aplica al caso LAP}

Este es, de los tres, el trabajo con la implicancia legal más favorable,
y conviene ser explícito sobre por qué.

El principio de minimización exige tratar solo los datos que resulten
necesarios para la finalidad. Un método que alcanza desempeño competitivo
con el 5\,\% de los datos etiquetados permite reducir en proporción
similar el volumen de imágenes del terminal que deben ser revisadas y
anotadas por personas. Menos imágenes anotadas significa menos exposición
de imágenes de pasajeros reales a anotadores humanos, menos copias del
material fuera del entorno controlado y una superficie de brecha
considerablemente menor. La eficiencia de etiquetado deja de ser una
ventaja de costo y se convierte en un argumento de cumplimiento.

Además, los experimentos de adaptación de dominio no supervisada del
artículo abordan de frente el problema que el proyecto tendrá que
resolver: ningún conjunto público de conteo de multitudes contiene
escenas de un terminal aéreo peruano.

\subsubsection*{Limitación reconocida por los autores}

Los autores documentan un caso de fallo atribuido al desenfoque producido
por el movimiento de la multitud, y lo presentan como tal dentro de su
comparación visual de predicciones. El método tiene, por tanto, un modo
de fallo reconocido y ligado a una condición, la multitud en movimiento,
que en un terminal aéreo es la norma y no la excepción. El artículo
no cuantifica en prosa la magnitud del error en ese caso concreto, de
modo que aquí se recoge la existencia del modo de fallo y no su tamaño.

Más importante aún, los propios autores registran la brecha de dominio.
Al transferir el modelo entre conjuntos sin adaptación de dominio el
desempeño se degrada de forma marcada, y son ellos quienes señalan
explícitamente que ese pobre desempeño subraya la brecha entre los datos
de origen y los de destino; cuando se incorporan al entrenamiento datos
no etiquetados del dominio de destino, los errores de estimación se
reducen de forma significativa. El artículo enuncia esa diferencia en
prosa sin consignarla numéricamente fuera de sus tablas.

Que los propios autores
adviertan una brecha de dominio de esta naturaleza implica que las
métricas publicadas sobre conjuntos públicos \textbf{no son
extrapolables} al Aeropuerto Jorge Chávez. Presentar a LAP una cifra de
exactitud tomada de la literatura, sin validación en el terminal, sería
una afirmación sin respaldo. Este informe recomienda, en consecuencia, que ninguna promesa cuantitativa de desempeño figure en un
entregable al cliente antes de una validación en sitio, aun cuando eso
implique presentar el proyecto sin cifras atractivas.

% ---------------------------------------------------------------------
\subsection{Síntesis: cómo se articulan los tres trabajos}
\label{sec:sintesis-antecedentes}

Los tres artículos no son alternativas entre las cuales elegir: forman
una progresión en la que cada uno resuelve una limitación del anterior y,
al hacerlo, desplaza el problema hacia el siguiente.

CSRNet~\cite{Li_2018_CVPR} resuelve el problema de \emph{cuántos} con una
arquitectura simple y entrenable de extremo a extremo, pero entrega una
superficie de densidad de la que no se puede extraer dónde está cada
persona. P2PNet~\cite{Song_2021_ICCV} resuelve el problema de
\emph{dónde} eliminando la representación intermedia y prediciendo puntos
directamente, pero paga ese avance con un costo de anotación que crece
con la densidad de la escena. P2R~\cite{Lin_2025_CVPR} resuelve el
problema del \emph{costo de anotación} llevando el enfoque por puntos al
régimen semi-supervisado, y de paso elimina el algoritmo húngaro que
P2PNet había introducido.

Vista desde el análisis de cumplimiento, la progresión tiene una lectura
distinta y en parte inversa. Cada paso hacia adelante en capacidad
técnica es también un paso hacia un perfil de riesgo distinto:

% 14/09/2026: longtable en lugar de tabularx dentro de center. Sin salto de
% pagina antes del Marco Teorico la tabla no cabia al pie de la pagina y
% dejaba un tercio en blanco; longtable puede partirse entre filas y repite
% la cabecera. Mismas columnas y mismo texto.
{\small
\setlength{\LTleft}{0pt}\setlength{\LTright}{0pt}
\begin{longtable}{@{}L{2.5cm}L{3.0cm}L{\dimexpr\linewidth-5.5cm-4\tabcolsep\relax}@{}}
\toprule
\textbf{Trabajo} & \textbf{Avance técnico} & \textbf{Desplazamiento del riesgo} \\
\midrule
\endfirsthead
\toprule
\textbf{Trabajo} & \textbf{Avance técnico} & \textbf{Desplazamiento del riesgo} \\
\midrule
\endhead
\bottomrule
\endlastfoot
CSRNet \cite{Li_2018_CVPR} &
Mapa de densidad sin representación individual &
Perfil de riesgo más bajo: la salida es agregada por construcción. El
riesgo se concentra aguas arriba, en qué se hace con el video original. \\
\addlinespace
P2PNet \cite{Song_2021_ICCV} &
Localización individual por puntos &
El riesgo se traslada a la salida del modelo: coordenadas por persona,
persistibles y encadenables en trayectorias. Exige control de
persistencia. \\
\addlinespace
P2R \cite{Lin_2025_CVPR} &
Desempeño competitivo con 5\,\% de etiquetas &
Reduce el riesgo del proceso de anotación (menos imágenes reales
expuestas a personas), pero introduce dependencia de pseudo-etiquetas
cuya calidad no es auditable caso por caso. \\
\end{longtable}
}

De esta lectura conjunta se desprenden tres criterios que el análisis de cumplimiento traslada al desarrollo
del sistema:

\begin{enumerate}[leftmargin=1.6em,itemsep=0.35em]
  \item \textbf{La granularidad de la salida es una decisión de
        cumplimiento, no solo de producto.} Si la finalidad declarada a
        LAP es estimar aglomeraciones, una salida agregada por zona
        satisface esa finalidad; una salida de coordenadas individuales
        la excede. El principio de proporcionalidad sugiere preferir la
        mínima granularidad que resuelva el problema, aunque el modelo
        sea capaz de más.
  \item \textbf{La eficiencia de etiquetado de P2R es un argumento de
        minimización aprovechable.} Conviene documentarla como tal en el
        análisis de riesgo y no solo como una ventaja de costo.
  \item \textbf{Ninguna cifra de la literatura es transferible al
        terminal sin validación local.} La propia evidencia de
        \cite{Lin_2025_CVPR} sobre brecha de dominio lo demuestra.
\end{enumerate}

Esta síntesis queda deliberadamente abierta en un punto. La elección de arquitectura del sistema no consta por escrito al cerrar esta
versión del informe, y por esa razón la auditoría de la
Sección~\ref{sec:auditoria} se formula de manera condicional sobre los
tres enfoques en lugar de sobre uno. Cuando esa elección quede
documentada, corresponde cerrar esta síntesis indicando cuál de los tres
se adopta y con qué granularidad de salida, porque de la granularidad
(mapa de densidad, puntos individuales o agregación por zona)
depende buena parte del examen de proporcionalidad de la
Sección~\ref{sec:proporcionalidad}.
